% !TEX root = CoeneDylanBP.tex
%%=============================================================================
%% Dataset en Dataverzameling
%%=============================================================================

\section{Dataset en Dataverzameling}
\label{sec:dataset}

Omdat de initiële opzet met een eigen cameraopstelling werd verlaten, werd overgeschakeld naar een aanpak op basis van bestaande wedstrijdbeelden. De dataset van \textcite{Lin2025}, die videomateriaal van competitieve trampolinesprongen bevat, werd gebruikt als primaire databron voor het trainen en evalueren van de modellen in de pipeline.

Deze dataset bestaat uit beelden van de Olympische Spelen en andere internationale wedstrijden, waarbij de camera niet volledig gefixeerd was maar meebewoog met de actie, vergelijkbaar met hoe een smartphone-applicatie in de praktijk zou worden ingezet.

Een direct gevolg van het werken met pre-opgenomen beelden is dat er geen controle bestond over de camerahoeken, instellingen of de positie ten opzichte van de trampoline.

Uit de dataset werden gelabelde afbeeldingen geselecteerd voor twee afzonderlijke trainingsdoeleinden:
\begin{itemize}
    \item Het trainen van het hoekpuntdetectiemodel (Sectie~\ref{sec:trampolinedetectie}): afbeeldingen van trampolines waarbij de vier hoeken van het bed handmatig werden geannoteerd.
    \item Het trainen van de landingsclassifier (Sectie~\ref{subsec:landingsmoment-classifier}): uitgeknipte beelden van het trampolinebed, voorzien van het label \textit{Landed} of \textit{Not Landed}, afhankelijk van of een springer in contact was met het doek.
\end{itemize}
